% !TeX root = ../教材题库_全部条目_转写工作稿.tex
% Source: XJ-B1-C06-PR-P250-Y0217
% 第6章 / 6.4 用样本估计总体 / 教材 p.250 / PDF 第257页
% OCR未校对：公式与文字均可能有误；未完成数学TeX转写；禁止直接发给学生
% 原题及相邻原解截图（核对用）：../题目原图/XJ-B1-C06-PR-P250-Y0217-01.webp
% 本文件只在转写完成后作为题干使用；OCR 条目尚不是可靠题干。
\begin{bankitem}{p.250 练习 第1题}{第6章 / 6.4 用样本估计总体 / 教材 p.250 / PDF 第257页}{XJ-B1-C06-PR-P250-Y0217}{OCR未校对：公式与文字均可能有误；未完成数学TeX转写；禁止直接发给学生}
\begin{ocrdraft}
[教材p.250 / PDF第257页]
\par
1.为了了解高一学生的体能情况，立德中学决定抽取部分学生进行1分钟跳
\par
绳测试，将所得数据整理后，得到一个频率分布直方图，如图所示．图中从左到
\par
右各小矩形面积之比为2：4：17：巧：9：3，第二小组频数为12．
\par
组距
\par
0
\par
90
\par
100110120130140150冫欠娄
\par
（第1题）
\par
(1)第二小组的频率是多少？样本容量是多少？
\par
（2）通过频率分布直方图估计总体的平均数、中位数、众数．
\end{ocrdraft}
\end{bankitem}
